\section{An unbounded positive forward-inverse kernel for binary--ternary products}
Let $q=(\sqrt5-1)/2$, $\phi=1+q=1/q$, and put
\[
 E(z)=(1+z)^u(1+z+z^2)^r,\quad d=u+2r,\qquad
 a_n=[z^n]\frac{E(z)}{1+qz}\quad(0\le n\le d).
\]
The sequence in the following theorem is extended by ZERO outside
$[0,d]$. The actual formal quotient has a nonzero alternating tail
after $d$; that tail is NOT asserted to be nonnegative or LC.

\begin{theorem}[All binary--ternary counts; a sharp pure-ternary threshold]
For EVERY $u\ge1,r\ge0$, the finite sequence $a_0,\ldots,a_d$
is strictly positive and globally log-concave. For $u=0$ it is
positive for every $r\ge0$ and is globally log-concave exactly when
$r=0$ or $r\ge6$.
\end{theorem}
\begin{proof}
Palindromicity and the defining forward recurrence give
\[
 \tau:=a_d=E(-q)=q^{2u}(2q^2)^r>0,\quad w=u+r,
 \quad a_{d-1}=\phi(1-\tau),\quad
 a_{d-2}=\phi\{w-\phi(1-\tau)\}.                 \tag{1}
\]
The last identity includes the case $d=1$, with $a_{-1}=0$,
since then $w=1$ and $E_{d-1}=E_0=1$. We only apply (1) when
$d\ge1$. In general (1) follows by solving the last two coefficient
equations backwards, using $E_d=1$ and $E_{d-1}=w$.

Suppose the old finite $a$ is positive and LC. Multiplication of E
by $1+z$ agrees with convolution of the zero-extended old prefix
at every rank through d; the only new coefficient at d+1 is
$(1-q)\tau=q^2\tau$, instead of the convolution's $\tau$.
It is positive and smaller. Convolution preserves LC, and lowering
this last coefficient improves the only affected interior minor.
This proves binary extension.

For multiplication by $1+z+z^2$, the coefficients through d again
agree with that positive LC convolution. Its last two coefficients
are changed to
\[
       a_{d-1}+q^2\tau,\qquad 2q^2\tau.           \tag{2}
\]
Both are positive; the coefficient at d+1 is lowered by $q\tau$,
so the minor at d improves. The ONLY new interior minor to pay is
at d+1. Substitution of (1) and $q^2+q=1$ gives its EXACT value
\[
 \{\phi(1-\tau)+q^2\tau\}^2
 -(\phi w-1+2\tau)2q^2\tau
       =\phi^2-2(\phi+qw)\tau.                    \tag{3}
\]
In particular an unbounded invariant sufficient for ternary extension
is $2(\phi+qw)\tau\le\phi^2$.

This invariant is paid at $(u,r)=(1,0)$: its slack is $4-5q>0$.
It is preserved by BOTH extensions. For the binary extension,
\[
 (\phi+qw)-q^2(\phi+q(w+1))=q(\phi+qw)-q^3\ge0.
\]
For the ternary extension and $w\ge1$,
\[
 (\phi+qw)-2q^2(\phi+q(w+1))
             =(2q-1)(q+qw-1)\ge0.                 \tag{4}
\]
Here $q>3/5$, so both factors on the right are positive.
Starting with $a=(1,q^2)$ and iterating binary and ternary extension
therefore proves ALL $u\ge1,r\ge0$, with no finite count bound.
The two-parameter potential induction itself is checked in Lean.

For pure ternaries use the SINGLE polynomial
\[
 (1+z+z^2)^6=(1,6,21,50,90,126,141,126,90,50,21,6,1).
\]
Its thirteen forward-inverse coefficients, eleven adjacent minors,
and the potential in (3) are explicitly positive by $q^2+q=1$.
The aggregate theorem in TernaryForwardInverseBase.lean verifies
ALL these displayed finite algebraic inequalities for the exact
positive quadratic root. The data file records the complete values,
not an unexplained search result. Equation (4) then proves every
$r\ge6$ by ordinary induction. For $1\le r\le5$, the penultimate
minors are respectively
\[
 1-\phi,\ 17-11\phi,\ 153-95\phi,\
 1137-703\phi,\ 7713-4767\phi,
\]
all strictly negative (direct comparison after squaring positive
rational sides). Thus six is sharp. Positivity under ternary
extension follows from (2) without any LC hypothesis. Starting at
the pure $r=1$ prefix $(1,q^2,2q^2)$ proves positivity for EVERY
pure count; the pure $r=0$ case is the one-term sequence $(1)$.
\end{proof}

The theorem supplies a finite positive LC convolution kernel at
EVERY count, including the diagonal product
$((1+z)(1+z+z^2)^2)^u$. It does not assert that an actual arm's
ordinary coefficients or positive differences equal this kernel.
In particular, it does not by itself provide an actual complete-C
entrance for the still unbounded diagonal $r=2u$.

\section{The complete positive companion for the actual diagonal recurrence}
Let $F_{u,r}=F_{1^\alpha2^u3^r}$ for the SAME $\alpha\ge1$,
$S_2=z+z^2$, $S_3=z+z^2+z^3$, $B_r=\sum_{j<r}S_3^j$,
and $K=S_2S_3^2$. Exact sorted-word counting gives
\[
 F_{u,r}=B_r+S_3^rF_{1^\alpha2^u},\qquad
 F_{1^\alpha2^{u+1}}=1+S_2F_{1^\alpha2^u}.
\]
Since $S_3-S_2=z^3$ and $(1-S_3)B_r=1-S_3^r$, it follows that
\[
 F_{u+1,r+2}=K F_{u,r}+P_r,\qquad
 P_r=1+S_3+S_3^2+z^3S_3^2B_r.                   \tag{5}
\]
The last positive summand is compulsory; a bare macro factor 233
does not give the actual arm recurrence.

For $H_3=(1-S_3)^{-1}$, let $E_{u,r}=H_3-F_{u,r}$.
Multiplying the geometric identity by $H_3$ gives
\[
 E_{u+1,r+2}=K E_{u,r}+z^3S_3^{r+2}H_3.          \tag{6}
\]
Put $d_3=(1+z)H_3$ and $L_{u,r}=(1-z)E_{u,r}/z^2$.
The two rank-zero/one losses vanish, so this is a legitimate series.
The complete shifted derivative-loss source is
\[
 L_{u+1,r+2}=K L_{u,r}+zS_3^{r+2}
                              +z^3S_3^{r+2}d_3. \tag{7}
\]
All coefficients of the two new source terms are nonnegative;
no part of the source has been silently discarded.
Three Lean commutative-ring identities verify (5)--(7) under
the explicitly stated word/geometric/renewal identities.

\paragraph{Scope and the remaining bridge.}
The new results are an unbounded finite-kernel LC theorem and exact
actual diagonal identities. The formal reports comprise eight
generic golden-root/potential/terminal-minor kernels, one aggregate
fixed degree-twelve base, and three complete diagonal algebra
kernels. Ordinary LC convolution, polynomial coefficient
identification, and the complete sequence induction above are
written proofs, not an end-to-end Lean formalization of compositions.
No count-bounded diagnostic is used as an unbounded premise.
The actual strong derivative-minor diagonal for $u\ge33$, general
earlier C, full A2 and ORIGINAL undilated coefficient unimodality
remain OPEN. The previously proved range $4\le u\le32,r\ge2u$
is not enlarged by this auxiliary kernel theorem alone.
